\documentclass[a4paper,11pt,fleqn]{article}%fleqn pour aligner les equations à gauche

\usepackage{../../Styles/Eval}

\newcommand{\chnum}{Chapitre 1}
\newcommand{\chtitle}{Composition initiale d'un système}
\newcommand{\dtype}{DS1}

\begin{document}
\maketitle

\section{On bulle ...}
\begin{enumerate}
    \item Calculer la quantité de matière $n_1$ de dioxyde de carbone contenue dans un volume $V = \SI{600}{mL}$ de gaz à \SI{20}{\celsius} et \SI{1013}{hPa}.
    \vspace{8em}
    \item La recharge d’un gazéificateur de boisson contient \SI{425}{g} de dioxyde de carbone dans un volume de \SI{600}{mL}. Calculer la quantité  de dioxyde de carbone contenu dans cette recharge.
    \vspace{8em}
    \item Comment est il possible que la recharge contienne une telle quantité de dioxyde de carbone ?
    \vspace{8em}
    \item Quand la recharge est considérée comme vide, que contient-elle ?
    \vspace{8em}
\end{enumerate}

\begin{figure}[b]
    \begin{minipage}{\linewidth}
        \caption*{\textbf{Données :}}
        Volume molaire des gaz à \SI{20}{\celsius} et \SI{1013}{hPa} : $V_m = \SI{24,0}{\liter\per\mole}$\\
Masse molaire du dioxyde de carbone : $M(\ce{CO2}) = \SI{44,0}{\gram\per\mole}$
    \end{minipage}
\end{figure}

\newpage
\section{Nettoyer à fond !}
L’eau de javel est une solution aqueuse contenant, entre autre, des ions hypochlorite \ce{ClO-}. Elle peut être commercialisée sous forme de bouteille ou de berlingot. La notice d’un berlingot contenant \SI{250}{mL} d’eau de javel indique de "verser le berlingot dans une bouteille d’un litre vide et compléter à l’eau froide".
\begin{enumerate}
    \item Calculer la concentration en quantité de matière des ions hypochlorite dans la solution préparée.
    \vspace{12em}
    \item Comparer la concentration à celle des ions hypochlorite contenus dans une bouteille commerciale d’eau de javel.
    \vspace{12em}
    \item Pour la désinfection des surface on recommande d’utiliser une solution à \SI{0,1}{\percent} d’ions hypochlorite, correspondant à une concentration de \SI{55E-3}{\mole\per\liter}. Proposer un protocole permettant de préparer \SI{500}{mL} de cette solution. Ce protocole précisera le matériel nécessaire ainsi que les calculs détaillés effectués.
    \vspace{12em}
\end{enumerate}
\begin{figure}[b]
    \begin{minipage}{\linewidth}
        \caption*{\textbf{Données :}}
        Concentration en ions hypochlorite dans un berlingot : \SI{0,46}{\mole\per\liter}\\
Concentration en ions hypochlorite dans une bouteille commerciale : \SI{0,11}{\mole\per\liter}
    \end{minipage}
\end{figure}\end{document}